\section{Part II: Jobs to Be Done: The Causal Anatomy of Demand}

\subsection{The Job Is the Progress, Not the Category}

The jobs-to-be-done framework developed by Clayton Christensen and colleagues emphasizes that customers do not simply buy products: they ``hire'' solutions in specific circumstances to make progress [R6, R7]. A commuter might purchase a morning milkshake not to savor confectionary craft, but to make a monotonous 45-minute highway drive engaging while keeping 1 hand on the wheel. In that specific circumstance, the milkshake competes with bagels, bananas, podcasts, and coffee, not solely with other milkshakes.

For a technical learner, ``take an AI course'' is a descriptive activity, not the underlying job. Possible causal jobs include:
\begin{itemize}[leftmargin=1.2em,itemsep=0.2ex]
  \item \textbf{Ownership capability:} Become capable of independently architecting, evaluating, and defending a production AI system.
  \item \textbf{Status defense:} Avoid professional embarrassment when technical leadership audits failure modes.
  \item \textbf{Role transition:} Credibly demonstrate systems engineering competence to qualify for a promotion.
  \item \textbf{Accountability structure:} Force oneself to maintain rigorous study habits when solitary motivation inevitably falters.
\end{itemize}

An educational institution may touch several of these jobs, but a specific buyer rarely values all of them equally. Diagnosing the dominant job governs every downstream choice: curriculum depth, instructional pacing, assessment rigor, pricing structure, and marketing message.

\subsection{Anatomy of an Actionable Job Statement}

An actionable job statement requires 3 mandatory components: a circumstance, a specific struggle, and desired progress under constraints.
\begin{center}
\small\itshape
When [specific circumstance], I struggle with [painful obstacle] and want to achieve [measurable progress], while respecting [unyielding constraints].
\end{center}

Weak: ``When engineers want AI education, they need interactive AI courses.'' This simply restates the founder's solution as the customer's need.

Strong: ``When backend engineers are assigned their first production AI feature, they struggle to evaluate hallucination rates, latency, and cost trade-offs; they want to make defensible architectural decisions before release review, while maintaining 40 weekly hours on their core backend duties.''

The second formulation illuminates the competitive set, project deadlines, evidentiary standards, and operational limits. It reveals whether the appropriate intervention is a multi-month course, an intensive 2-day architectural review studio, an enterprise evaluation harness, or a specialized diagnostic tool. The correct product architecture remains an open question until the job is crystallized.

\subsection{Functional, Emotional, and Social Dimensions}

A functional job specifies physical or technical progress: ``deploy a reliable semantic search service that maintains p99 latency under 200 milliseconds.'' Emotional progress specifies personal psychological relief: ``stop feeling out of depth when executive leadership questions model hallucinations.'' Social progress specifies external perception: ``be recognized by senior engineering peers as a rigorous, dependable systems architect.''

To uncover these intertwined dimensions during discovery, ask about \textbf{consequences and reactions}:
\begin{itemize}[leftmargin=1.2em,itemsep=0.2ex]
  \item What did success or failure change in your weekly standing at work?
  \item Who noticed when the system crashed, and what was said?
  \item What was the most professionally uncomfortable moment of the project?
  \item What were you desperate to avoid having to explain to management?
\end{itemize}
Respect personal boundaries; never manufacture artificial vulnerability. Emotional and social jobs do not suggest irrationality: they reflect the genuine human stakes of organizational life.

\subsection{The 4 Forces Governing Progress}

A customer's decision to adopt a new solution is governed by 4 interacting forces:
\begin{enumerate}[leftmargin=1.2em,itemsep=0.2ex]
  \item \textbf{Push of the current situation:} Frustration with current errors, rising costs, mounting deadlines, and the failure of existing workarounds.
  \item \textbf{Pull of the new solution:} The magnetic appeal of promised capabilities, superior outcomes, and professional mastery.
  \item \textbf{Anxiety of the new solution:} Fear of hidden complexity, wasted money, public failure, demanding workloads, or an untrusted provider.
  \item \textbf{Habit and allegiance to the present:} Familiar comfort with established spreadsheets, existing codebases, trusted colleagues, and default routines.
\end{enumerate}

A commercial transaction occurs only when the forces favoring change (\textbf{Push + Pull}) substantially overpower the forces resisting change (\textbf{Anxiety + Habit}).

Most novice founders pour all their effort into magnifying Pull through louder marketing and grander promises. Yet adoption is far more frequently blocked by \textbf{Anxiety} (fear of wasting scarce money on an unproven credential) and \textbf{Habit} (comfort with free YouTube tutorials). Reducing anxiety through transparent syllabi, small diagnostic trials, and explicit prerequisite audits is often 10 times more effective than doubling promotional claims.
